% =========================================================
% mechlib/elements/trusses.tex
% FACHWERKE: STAB UND GEMEINSAME GELENKDARSTELLUNG
% =========================================================
% Dieses Modul enthaelt die atomaren Fachwerkelemente. Die Topologie eines
% Fachwerks bleibt bewusst im Anwendungsdokument: Knoten werden als TikZ-
% Koordinaten definiert und durch mech/rod bzw. \mechJoints verbunden.
%
% Oeffentliche API:
%   mech/rod     Stab zwischen "from" und "to" mit optionalem Label
%   \mechJoints{A,B,C,...} zeichnet alle angegebenen Gelenke
%
% Darstellungsregeln:
%   - Staebe verwenden standardmaessig den Stil mech/rod: schwarz und duenn.
%   - Staebe werden auf dem Layer "mech background" gezeichnet.
%   - Stabnummern besitzen keine Fuellung (fill=none) und bleiben transparent.
%   - Gelenke liegen optisch vor den Staeben und sind mit mechBodyColor gefuellt.
%
% Beispiel:
%   \pic {mech/rod={from=(A),to=(B),label=1,label side=above}};
%   \mechJoints{A,B}
% =========================================================
% -------------------------------------------------------------------------
% Einzelner Fachwerkstab
% -------------------------------------------------------------------------
% "style" erlaubt lokale Abweichungen, ohne den globalen mech/rod-Stil zu
% veraendern. Die Standardbeschriftung liegt an der Stabmitte.
\pgfkeys{
  /mech/rod/.is family,
  /mech/rod,
  from/.initial={(0,0)},
  to/.initial={(1,0)},
  label/.initial=,
  label side/.initial=above,
  label style/.initial={mech/label,fill=none,text=black},
  style/.initial=mech/rod,
}
\tikzset{
  pics/mech/rod/.style args={#1}{
    code={
      \pgfkeys{/mech/rod/.cd,#1}
      \edef\mechRodFrom{\pgfkeysvalueof{/mech/rod/from}}
      \edef\mechRodTo{\pgfkeysvalueof{/mech/rod/to}}
      \coordinate (-start) at \mechRodFrom;
      \coordinate (-end) at \mechRodTo;
      \coordinate (-center) at ($(-start)!.5!(-end)$);
      \begin{pgfonlayer}{mech background}
        \draw[\pgfkeysvalueof{/mech/rod/style}] (-start)--(-end);
      \end{pgfonlayer}
      \node[\pgfkeysvalueof{/mech/rod/label style},\pgfkeysvalueof{/mech/rod/label side}]
        at (-center) {\pgfkeysvalueof{/mech/rod/label}};
    }
  },
  pics/mech/rod/.default={},
}
% -------------------------------------------------------------------------
% Gelenkliste
% -------------------------------------------------------------------------
% Das Makro erwartet eine kommaseparierte Liste bereits definierter TikZ-
% Koordinaten. Jedes Gelenk wird identisch gezeichnet. Die Gelenke sollten nach
% den Staeben aufgerufen werden, damit sie die Stabenden optisch ueberdecken.
\newcommand{\mechJoints}[1]{%
  \foreach \mechJoint in {#1}{%
    \draw[fill=mechBodyColor,draw=black,line width=.2mm]
      (\mechJoint) circle[radius=.6mm];%
  }%
}
